
% \documentclass[a4paper,12pt]{article}
\documentclass[uplatex]{jsarticle}
\usepackage{otf}

\usepackage{amsmath, amssymb, bm}
\usepackage{geometry}
\usepackage{hyperref}
\geometry{margin=25mm}
\title{オペレータモデルを含むエネルギー保存則によるハプティクスフィードバック生成法}
\author{}
\date{}

\begin{document}
\maketitle

% \section{Overview}

\section{テレオペレーションシステムにおけるエネルギー収支}
ロボットとデバイスからなるシステムのエネルギーの総量を考える。
ロボットが外界と、デバイスがオペレータと、それぞれ力を及ぼし合うことで、以下のようにエネルギーが流れる。
\begin{itemize}
\item ロボットの手先が外界に接触し、システムにエネルギーが流入する。
\begin{equation}
P_{in}^{env}(t) = \bm{w}_{env}^T(t) \bm{v}_{robot}(t)
\end{equation}
\item デバイスがオペレータに力を提示し、手が動くことでシステムからエネルギーが流出する。
\begin{equation}
P_{out}^{op}(t) = \bm{w}_{fb}^T(t) \bm{v}_{hand}(t)
\end{equation}
\item オペレータが提示力に抗い、手を動かし返すことで、システムにエネルギーが流入する。
\begin{equation}
P_{in}^{op}(t) = \bm{w}_{hand}^T(t) \bm{v}_{hand}(t)
\end{equation}
\item ロボットの手先が外界に力を加えることで、システムからエネルギーが流出する。
\begin{equation}
P_{out}^{env}(t) = \bm{w}_{robot}^T(t) \bm{v}_{robot}(t)
\end{equation}
\end{itemize}
ただし$\bm{w} \in \mathbb{R}^6$はレンチ、$\bm{v} \in \mathbb{R}^6$はツイストである。

予測されるエネルギー変化は
\begin{equation}
\Delta E^{pred} = (P_{out}^{op}(t) + P_{out}^{env}(t) - P_{in}^{env}(t) - P_{in}^{op}(t)) \Delta t
\end{equation}

ここで、
\begin{itemize}
\item $\Delta E^{pred} > 0$: システムのエネルギーは減少 (安全)
\item $\Delta E^{pred} < 0$: システムのエネルギーは増加 (危険)
\end{itemize}

以下でシステムから出入りするエネルギー量を求め、系のエネルギーが増加しないようなフィードバックレンチの決定法を考える。
以下では各時間ステップごとに離散化して考える。
ステップのインデックスを$i$とする。


\section{外界からシステムに流入するエネルギー}
ステップ$i$のときロボットが外界から受けるレンチ$\bm{w}_{measured}(i)$は、ロボットにつけられたFTセンサから計測される。
またステップ$i$でのロボットの手先のツイスト$\bm{v}_{robot}(i)$は、各関節の角速度からFKを解くことで計測できる。
よってステップ$i$からステップ$i+1$の間にロボットの手先が外界から力をうけることでシステムに流入するパワーは以下のように求められる。
\begin{equation}
P_{in}^{env} = \bm{w}_{measured}^T(i) \bm{v}_{robot}(i)
\end{equation}


\section{システムからオペレータに流出するエネルギー}

\subsection{オペレータのモデル化}
オペレータのインピーダンスモデルは以下で与えられる。
\begin{equation}\label{operator_model_eq}
F(t) = M \ddot{x}(t) + B \dot{x}(t) + K (x(t) - x_{ref})
\end{equation}
ここで、$F(t)$はオペレータの手に与えられる力で、$M, B, K$ はそれぞれ等価質量、粘性、剛性であり、時間に依存して変化する。
また$x_{ref}$はオペレータの参照原点であり、テレオペ開始時の脱力した状態での位置、すなわちテレオペの指令送信で用いるオペレータ系の初期位置を用いる。
このモデルが6次元のレンチの各要素についてそれぞれ成立すると仮定する。

離散時間表現により、次の線形回帰式に変換できる。
\begin{equation}
F_i \approx \bm{\phi}_i^T \bm{\theta}_i \quad 
\bm{\phi}_i = \begin{bmatrix} \ddot{x}_i & \dot{x}_i & x_i - x_{ref} \end{bmatrix}^T
\quad \bm{\theta}_i = \begin{bmatrix} M_i & B_i & K_i \end{bmatrix}^T
\end{equation}

\subsection{オペレータのインピーダンス推定}
先の離散時間表現は最小二乗法に適しており、逐次最小二乗法(RLS)によってオンライン推定できる。
オペレータのインピーダンス$\bm{\theta}_i = \begin{bmatrix} M_i & B_i & K_i \end{bmatrix}^T$を$\bm{\theta}_{i-1}$から推定する。
RLSによる更新式は以下の通り(詳しくは\url{https://manabitimes.jp/math/2084}を参照)
\begin{equation}
Q_i = \frac{1}{\lambda} (Q_{i-1} - \frac{Q_{i-1} \bm{\phi}_i \bm{\phi}_i^T Q_{i-1}}{\lambda + \bm{\phi}_i^T Q_{i-1} \bm{\phi}_i})
\end{equation}
\begin{equation}
\bm{\theta}_i = \bm{\theta}_{i-1} + Q_i \bm{\phi}_i (y_i - \bm{\phi}_i^T \bm{\theta}_{i-1})
\end{equation}
ただし$\bm{\phi}_i = \begin{bmatrix} \ddot{x}_i & \dot{x}_i & {x}_i \end{bmatrix}^T$で、$y_i = f_{feedback}(i-1)$は前ステップにオペレータに提示したレンチ成分である。

\subsection{システムから流出するエネルギーの予測}
ロボットで計測されたレンチをスケールしたもの$\bm{w}^{nom}(i) = \bm{k} \odot \bm{w}_{robot}^{measured}(i)$をオペレータに提示した際のオペレータの即時応答速度を予測し、システムに流入するエネルギーを予測する。
($\odot$はアダマール積で、各要素ごとの積である)

式(\ref{operator_model_eq})よりオペレータの手の加速度は以下のように求められる。
(位置、姿勢の各要素について)
\begin{equation}
\ddot{x}(i) = \frac{1}{M_i} (F^{nom}(i) - B_i \dot{x}(i) - K_i x(i))
\end{equation}
これよりオペレータの手の応答速度は以下のように予測できる。
\begin{equation}
\dot{x}(i+1) \approx \dot{x}(i) + \ddot{x}(i) \Delta t = \dot{x}(i) + \frac{\Delta t}{M_i} (F^{nom}(i) - B_i \dot{x}(i) - K_i x(i))
\end{equation}
これを6次元それぞれで予測することで、レンチを提示した際のオペレータの手の応答ツイスト$\bm{v}_{hand}^{pred}$が予測できる。

これよりシステムから流出するパワーは以下のように予測できる。
\begin{equation}
P_{out}^{op} = (\bm{w}^{nom}(i))^T \bm{v}_{hand}^{pred}(i+1)
\end{equation}


\section{オペレータからシステムに流入するエネルギー}
提示力を受けたオペレータは、それに対して押し返すように手を動かす。
オペレータが発揮するレンチは、先のオペレータのモデルを用いて推定する。
ステップ$i$でのオペレータの手の位置、速度、加速度はMocapから測定可能である。
(ただしノイズの考慮が必要)
この測定値からオペレータの発揮レンチの各成分は以下のように推定できる。
\begin{equation}
F_{op,est}(i) = M_i \ddot{x}(i) + B_i \dot{x}(i) + K_i x(i)
\end{equation}
またオペレータのツイスト$\bm{v}_{op}(i)$はMocapから計測できる。
よってオペレータからシステムに流入するパワーは以下のように表せる。
\begin{equation}
P_{in}^{op} = \bm{w}_{op,est}^T(i) \bm{v}_{op}(i)
\end{equation}


\section{システムから外界に流出するエネルギー}
ステップ$i$のときロボットが外界に及ぼすレンチ$- \bm{w}_{measured}(i)$は、ロボットにつけられたFTセンサでの計測値の反力として計測できる。
またステップ$i$でのロボットの手先のツイスト$\bm{v}_{robot}(i)$は、各関節の角速度からFKを解くことで計測できる。
よってステップ$i$からステップ$i+1$の間にロボットの手先が外界にレンチを及ぼすことでシステムに流入するパワーは以下のように求められる。
\begin{equation}
P_{out}^{env} = \bm{w}_{measured}^T(i) \bm{v}_{robot}(i) = P_{in}^{env}
\end{equation}


\section{システムのエネルギー収支と受動性保持のためのフィードバックレンチの決定法}
以上までの議論から１サイクルでのエネルギー収支は以下のようになる。
\begin{equation}
\begin{split}
\Delta E^{pred} &= (P_{out}^{op}(t) + P_{out}^{env}(t) - P_{in}^{env}(t) - P_{in}^{op}(t)) \Delta t \\
&= ((\bm{w}^{nom}(i))^T \bm{v}_{hand}^{pred}(i+1) - \bm{w}_{op,est}^T(i) \bm{v}_{op}(i)) \Delta t
\end{split}
\end{equation}

\subsection{エネルギータンクによる受動性保持}
システムの保有するエネルギーとして、$E_{tank}$を定義する。
1章での議論の通り、$\Delta E^{pred} > 0$のときシステムからエネルギーが流出する。
このとき$\Delta E^{pred} > E_{tank}$となると、システムに蓄えられたエネルギー以上をオペレータに流出させるために、システムが新たにエネルギーを生み出すことになり受動性を破る。
そのため$\Delta E^{pred} <= E_{tank}$となるように流出エネルギーを減少させる必要がある。
以下のようにフィードバックレンチを減少させる。

$\Delta E^{pred} > E_{tank}$のとき
\begin{equation}
\alpha = \text{min} (1,\frac{E_{tank}}{\Delta E^{pred} + \epsilon}), \alpha \in [0,1)
\end{equation}
\begin{equation}
\bm{w}_{cand} = \alpha \bm{w}^{nom}
\end{equation}

$0 < \Delta E^{pred} < E_{tank}$のときは蓄えられたエネルギーを流出させる。
\begin{equation}
E_{tank} -= \Delta E^{pred}
\end{equation}
このときのフィードバックレンチは
\begin{equation}
\bm{w}_{cand} = \bm{w}^{nom}
\end{equation}
  

\subsection{TDPAによる受動性保持}
1章での議論の通り、$\Delta E^{pred} < 0$のときシステムのエネルギーは増加し、受動性が破れるおそれがあるため危険である。
このようにシステムがエネルギーを自発的に生み出すことを避けるために、TDPA(Time Domein Passivity Approach)による受動性の保持を行う。
PO(Passivity Observer)により、これまでシステムから外部に流出した累積エネルギー$E_{PO}$を監視する。
\begin{equation}
\begin{split}
E_{PO}(i+1) &= E_{PO}(i) + \Delta E_{PO}(i) \\ 
&= E_{PO}(i) + (\bm{w}_{cand}^T(i) \bm{v}_{hand}^{pred}(i) - \bm{w}_{op,est}^T(i) \bm{v}_{op}(i)) \Delta t
\end{split}
\end{equation}
$E_{PO} < 0$のとき、システムがエネルギーを生み出すことになるため、危険である。
よってエネルギーを散逸させて受動性を保持するために、PC(Passivity Controller)によりフィードバックレンチにダンピングを追加する。
\begin{equation}
\bm{w}_{feedback}(i) = \bm{w}_{cand}(i) - D_{add} \bm{v}_{hand}^{pred}(i)
\end{equation}
このようにダンパーを追加することによって、システムからエネルギーを散逸させ、自発的にエネルギーを生み出すことを避ける。
散逸したエネルギーはエネルギータンクに回収する。(エネルギータンクの上限以内)

$D_{add} \in \mathbb{R}^{6 \times 6}$は次のように求める。
ダンピングを入れたあとに系から流出するパワーは
\begin{equation}
\begin{split}
\frac{\Delta E_{PO}(i)}{\Delta t} &= \bm{w}_{feedback}^T(i) \bm{v}_{hand}^{pred}(i) - \bm{w}_{op,est}^T(i) \bm{v}_{op}(i) \\
&= (\bm{w}_{cand}(i) - D_{add} \bm{v}_{hand}^{pred}(i))^T \bm{v}_{hand}^{pred}(i) - \bm{w}_{op,est}^T(i) \bm{v}_{op}(i) \\
&= \bm{w}_{cand}^T(i) \bm{v}_{hand}^{pred}(i) - \bm{v}_{hand}^{pred,T}(i) D_{add} \bm{v}_{hand}^{pred}(i) - \bm{w}_{op,est}^T(i) \bm{v}_{op}(i)
\end{split}
\end{equation}
システムの受動性を守るための条件は、システムがエネルギーを生み出さないこと、すなわち流出エネルギーが0以上であることなので
\begin{equation}
E_{PO}(i+1) = E_{PO}(i) + \Delta E_{PO}(i) \geq 0
\end{equation}
代入して整理すると
\begin{equation}
\bm{v}_{hand}^{pred,T}(i) D_{add} \bm{v}_{hand}^{pred}(i) &\leq \bm{w}_{cand}^T(i) \bm{v}_{hand}^{pred}(i) - \bm{w}_{op,est}^T(i) \bm{v}_{op}(i) + \frac{E_{PO}(i)}{\Delta t} := P_{req}
\end{equation}
以上の関係を満たす$D_{add}$を求める。

ダンピング係数$D_{add}$は各軸ごとに要素を持つ対角行列として決定する。
\begin{equation}
D_{add} = diag(d_1, \ldots, d_6) (d_i > 0)
\end{equation}





\section{Implementation as ROS Node}
ROSノードは以下の入出力を持つ：
\begin{itemize}
\item Subscriber:
  \begin{itemize}
    \item /cfs/data (WrenchStamped) : ロボットの手先で計測されたレンチ
    \item /me6\_robot/joint\_states (JointStates) : ロボットの関節角度
    \item /twin\_hammer/mocap/pose (PoseStamped) : デバイス位置姿勢
  \end{itemize}
\item Publisher:
  \begin{itemize}
    \item /twin\_hammer/haptics\_wrench (WrenchStamped) : フィードバックレンチ
  \end{itemize}
\end{itemize}

ノード内の処理ループ：
\begin{enumerate}
\item オペレータのインピーダンスモデルをRLSにより逐次推定
\item 計測レンチをスケーリングし，オペレータ応答を予測
\item $\Delta E^{pred}$ を計算し，タンク残量と比較
\item 受動性を守るため，必要に応じてスケーリング係数$\alpha$や追加ダンピング$B_d$を調整
\item フィードバックレンチを生成しパブリッシュ
\end{enumerate}

\section{Tuning Parameters}
チューニングが必要となる主要パラメータは以下：
\begin{itemize}
\item スケーリング係数 $\alpha$: 値を大きくすると感触は直感的だが，不安定になりやすい．小さくすると安定だが感触が弱い．
\item ダンピング係数 $B_d$: 増加させると安定性は高まるが，操作感が重くなる．
\item エネルギータンク初期値 $E_{tank}(0)$: 小さすぎるとすぐ介入が発生，大きすぎると不安定化を許す可能性がある．
\end{itemize}

\end{document}
